% Part: methods
% Chapter: induction
% Section: strong-induction

\documentclass[../../../include/open-logic-section]{subfiles}

\begin{document}

\olfileid{mth}{ind}{str}

\olsection{బలమైన ఆగమనం}

పైన చెప్పిన ఆగమన సూత్రంలో $P(0)$ను, అలాగే $P(k)$ అయితే $P(k+1)$ను నిరూపిస్తాం. రెండవ భాగంలో $P(k)$ సత్యమని తీసుకొని, దానితో $P(k+1)$ను నిరూపిస్తాం. సమానంగా, ధన సూచికల సందర్భంలో $P(k-1)$ను తీసుకొని దానితో $P(k)$ను నిరూపించవచ్చు---ముఖ్యమైనది ఏ సంఖ్య నుంచి దాని అనువర్తి సంఖ్యకు నిగమనాన్ని చేయగలగడం; పూర్వ సంఖ్యకు వాదన సత్యమని తీసుకొని ఆ సంఖ్యకు వాదనను నిరూపించగలగడం.
\sourcecorrection{OLTEMTHINDSTR-001}{మూలం $P(k-1)$ రూపాన్ని ఎలాంటి పరిమితి లేకుండా సమానమైనదిగా చెప్పింది; సహజ సంఖ్యలలో శూన్యానికి పూర్వ సంఖ్య లేదు. ధన సూచికలకే ఆ పునర్వ్యాఖ్యను పరిమితం చేశాం.}

ఆగమన సూత్రానికి మరో రూపం ఉంది. ఇందులో $k$కు ముందరి సంఖ్య $k-1$కు మాత్రమే వాదన సత్యమని అనుకోకుండా,~$k$ కంటే చిన్న అన్ని సంఖ్యలకూ సత్యమని తీసుకొని, దానితో~$k$కు స్థాపిస్తాం. దీనివల్ల కూడా అన్ని~$n \in
\Nat$లకు $P(n)$ వస్తుంది. ఎందుకంటే $P(0)$ను స్థాపించిన వెంటనే~$1$ కంటే చిన్న అన్ని సంఖ్యలకు $P$ సత్యమని స్థాపించినట్లే. అన్ని $l<k$లకు $P(l)$ అయితే $P(k)$ అని తెలిస్తే, ప్రత్యేకంగా $k=1$ సందర్భానికీ అది వర్తిస్తుంది. కాబట్టి $P(1)$ను తేల్చవచ్చు. అప్పుడు $P(0)$, $P(1)$ రెండూ నిరూపించాం, అంటే అన్ని $l<2$లకు $P(l)$; అలాగే అన్ని $l<2$లకు $P(l)$ అయితే $P(2)$ అనే సోపాధికం కూడా ఉంది కనుక~$P(2)$ను తేల్చవచ్చు; ఇలా కొనసాగుతుంది.

నిజానికి ``అన్ని $k$లకు, అన్ని $l<k$లకు $P(l)$ అయితే $P(k)$'' అనే సాధారణ సోపాధికాన్ని స్థాపించగలిగితే, $P(0)$ను విడిగా స్థాపించనవసరం లేదు: అది ఆ సోపాధికం నుంచే వస్తుంది. ``అన్ని $l<k$లకు $P(l)$'' వంటి సార్వత్రిక వాదన, $l<k$ అయ్యేవి ఏవీ లేకపోతే సత్యమని గుర్తుంచుకోండి. ఇది ఖాళీ పరిధిపై పరిమాణీకరణ: ``అన్ని $A$s, $B$s'' అనేది $A$s ఏవీ లేకపోతే సత్యం; $x$ ఏదీ~$!A(x)$ను తృప్తిపరచకపోతే $\lforall[x][(!A(x) \lif !B(x))]$ సత్యం. ఇక్కడ ఆకారిక రూపం ``$\lforall[l][(l < k \lif P(l))]$''; $l < k$ అయ్యేవి లేకపోతే అది సత్యం. $k=0$ అయితే అదే పరిస్థితి: $l<0$ అయ్యేది ఏదీ లేదు; కాబట్టి ``అన్ని $l<0$లకు $P(l)$'' అనేది $P$ ఏదైనా సత్యం. అందువల్ల ``అన్ని $l<k$లకు $P(l)$ అయితే $P(k)$'' నిరూపణ స్వయంగా~$P(0)$ను స్థాపిస్తుంది.
\sourcecorrection{OLTEMTHINDSTR-002}{మూలం శూన్య సందర్భంలో ``అన్ని $l<0$లకు $P(0)$'' అని రాసింది; ఖాళీ పరిధిపై అది కూడా సత్యమే అయినా, ఇక్కడి సాధారణ ఆగమన పూర్వపక్షం ``అన్ని $l<0$లకు $P(l)$''. దానికి అనుగుణంగా సూచికను సరిచేశాం.}

ఈ రూపం,~$k$కు వాదనను నిరూపించడానికి $k-1$ సందర్భం ఒక్కటే సరిపోక, ఒకటి లేదా అంతకంటే ఎక్కువ $l<k$ సందర్భాల పరికల్పన అవసరమయ్యే చోట ఉపయోగపడుతుంది.

\end{document}
